% AIbuliding 訓練資料草稿格式（給人工寫解答用）
% 用法：複製本檔 → 依題填寫 → 交給 agent 轉成 JSONL
% 規則：先寫進 data/draft.jsonl；審過再 promote 到 data/train.jsonl
%
% 標記說明（請保持關鍵字一致，方便之後自動匯入）：
%   \qa{id}{...}     一題一塊
%   \question{...}   題目（user）
%   \solution{...}   逐步解答（assistant）
%   解答內請用：Step 1： / Step 2： / **答案：** / **常見錯誤：** 或 **注意：**
%
% 有圖的題：把圖中條件寫進 \question{}，例如：
%   [圖] 長方形 $ABCD$，其中 $AB=6$，$AD=4$，對角線交於 $O$。

\documentclass{article}
\usepackage{amsmath,amssymb}
\usepackage[utf8]{inputenc}
\begin{document}

% ========== 題目 1 ==========
\qa{2015p1_q01}{
\question{
化簡 $\dfrac{m^{9}}{(m^{3}n^{7})^{5}}$，並以正指數表示答案。
}
\solution{
Step 1：先處理括號內的指數。

$$\left(m^{3}n^{7}\right)^{5}=m^{15}n^{35}$$

Step 2：相除時同底數指數相減：

$$\dfrac{m^{9}}{m^{15}n^{35}}=m^{9-15}n^{-35}=m^{-6}n^{-35}$$

Step 3：改寫成正指數：

$$\dfrac{1}{m^{6}n^{35}}$$

**答案：** $\dfrac{1}{m^{6}n^{35}}$

**常見錯誤：** 忘記外層指數同時作用在 $m$ 與 $n$。
}
}

% ========== 題目 2（有圖示例）==========
\qa{2015p1_q08}{
\question{
在圖中，長方形 $ABCD$ 中，$AB=6$，$AD=4$，對角線 $AC$ 與 $BD$ 交於 $O$。
求 $AO$ 的長度。

[圖] 長方形 $ABCD$，$AB=6$，$AD=4$，對角線交於 $O$。
}
\solution{
Step 1：長方形對角線相等且互相平分，故 $O$ 為 $AC$ 中點。

Step 2：由畢氏定理，

$$AC=\sqrt{AB^{2}+AD^{2}}=\sqrt{6^{2}+4^{2}}=\sqrt{52}=2\sqrt{13}$$

Step 3：因此

$$AO=\dfrac{1}{2}AC=\sqrt{13}$$

**答案：** $\sqrt{13}$

**注意：** 長方形對角線互相平分，不一定互相垂直。
}
}

\end{document}
